% 核心程序片段（脚本生成；完整工程见支撑材料 src/）
\lstset{language=Python,basicstyle=\ttfamily\scriptsize,numbers=left,
  numberstyle=\tiny,breaklines=true,breakatwhitespace=false,
  columns=flexible,keepspaces=true,showstringspaces=false,tabsize=4,frame=none}
\section{核心程序片段}
以下给出论文方法对应的核心算法片段（线性规划矩阵构建、凸组合权重标定、裕度参数确认、逐槽因果执行）；完整可运行工程与运行环境见支撑材料 \texttt{src/}（支持 uv 与 requirements.txt 两种安装方式）。
\subsection*{一、线性规划矩阵构建}
\noindent\texttt{program/src/solve/core/lp.py :: plan\_lp}\par
\lstinputlisting[firstline=53,lastline=117]{../program/src/solve/core/lp.py}
\noindent\texttt{program/src/solve/core/lp.py :: adjust\_day}\par
\lstinputlisting[firstline=240,lastline=308]{../program/src/solve/core/lp.py}
\subsection*{二、凸组合预测与参数标定}
\noindent\texttt{program/src/solve/models/weights.py :: softmax}\par
\lstinputlisting[firstline=17,lastline=20]{../program/src/solve/models/weights.py}
\noindent\texttt{program/src/solve/models/weights.py :: simplex\_grid}\par
\lstinputlisting[firstline=23,lastline=32]{../program/src/solve/models/weights.py}
\noindent\texttt{program/src/solve/consistency.py :: ewma\_weights\_from\_table}\par
\lstinputlisting[firstline=67,lastline=85]{../program/src/solve/consistency.py}
\noindent\texttt{program/src/solve/q2\_tune.py :: seqs\_ewma}\par
\lstinputlisting[firstline=46,lastline=53]{../program/src/solve/q2_tune.py}
\noindent\texttt{program/src/solve/models/adaptive.py :: day\_cost}\par
\lstinputlisting[firstline=269,lastline=280]{../program/src/solve/models/adaptive.py}
\noindent\texttt{program/src/solve/models/adaptive.py :: \_one\_day}\par
\lstinputlisting[firstline=329,lastline=349]{../program/src/solve/models/adaptive.py}
\noindent\texttt{program/src/solve/models/adaptive.py :: calibrate}\par
\lstinputlisting[firstline=299,lastline=324]{../program/src/solve/models/adaptive.py}
\subsection*{三、裕度参数确认}
\noindent\texttt{program/src/solve/consistency.py :: kappa\_load}\par
\lstinputlisting[firstline=108,lastline=110]{../program/src/solve/consistency.py}
\noindent\texttt{program/src/solve/consistency.py :: margin\_pv}\par
\lstinputlisting[firstline=113,lastline=115]{../program/src/solve/consistency.py}
\noindent\texttt{program/src/solve/q2\_tune.py :: margin\_refs}\par
\lstinputlisting[firstline=100,lastline=122]{../program/src/solve/q2_tune.py}
\noindent\texttt{program/src/solve/q2\_tune.py :: rolling\_kappa}\par
\lstinputlisting[firstline=125,lastline=134]{../program/src/solve/q2_tune.py}
\subsection*{四、逐槽因果执行}
\noindent\texttt{program/src/solve/core/causal.py :: exec\_segment\_causal}\par
\lstinputlisting[firstline=44,lastline=118]{../program/src/solve/core/causal.py}
